\section{Limitations}
\label{sec:limitations}

\emph{One-dimensional inputs.} The geometry experiments and
\cref{prop:length,prop:sine,prop:floor} are for curves in the plane; claim (G)
has not been measured for surfaces.
\emph{A degenerate family.} For fixed $w$ the functions $p\sin(w\pi x+\varphi)$
span a two-dimensional space, and their covariance kernel
$\tfrac{p^2}{2}\cos(w\pi(x-x'))$ has rank 2. No member can fit more than two
generic points at once, at any frequency, and the mean per-point distance
cannot see this because it scores each point alone.
\emph{Nothing to learn.} The labels are independent of the inputs, so the
experiments concern the loss at initialisation only and say nothing about
learning speed or generalisation.
\emph{Expectation only.} \Cref{thm:marginal} and \cref{cor:he} are about the
mean initial loss; whether the probability of a low initial loss changes with
width is not addressed.
\emph{Own scripts.} The experiments are small computations written for this
question, not an established benchmark, and the code has needed correction more
than once (\cref{app:wrong}).
\emph{Proof status.} The constant in \cref{prop:sine} rests on a leading-order
argument and a numerical check.

\section{Discussion}
\label{sec:discussion}

The conjecture joined two facts that are each true: a squigglier curve is
nearer to random points, and larger networks learn faster. The link fails at
both ends. Width does not make a random network's curve longer, and the
nearness that squiggliness buys is nearness to points displaced in the input
direction, which a regression loss does not reward and which
\cref{thm:sandwich} shows is paid for in slope.

\paragraph{Where ``starts closer'' might still hold.} The claim could be moved
to weight space: a larger network may need to travel a shorter relative distance
from its initial weights to a fit. The repository's notes predict from the
lazy-training results of \cref{sec:related} that it does, in which case an
experiment would be a replication. We have not run it.

\paragraph{An untested proposal: joint fit.} The nearest-point distance scores
each point alone. A cost that sees all points at once is $y^\top K^{-1}y$, where
$y\in\R^n$ collects the labels and $K$ is the positive-definite covariance
matrix of the random function at the $n$ inputs; it is the data-fit term of the
Gaussian-process marginal likelihood \citep{rasmussen2006gaussian}. For
uncorrelated labels of zero mean and unit variance its mean is $\tr(K^{-1})$,
and $\tr(K^{-1})\ge n^2/\tr(K)$ with equality if and only if $K$ is a multiple
of the identity, by the Cauchy--Schwarz inequality
$n^2=\bigl(\sum_i\lambda_i^{1/2}\lambda_i^{-1/2}\bigr)^2\le\sum_i\lambda_i\sum_i\lambda_i^{-1}$
for the eigenvalues $\lambda_i$ of $K$. At fixed output variance $\tr(K)$ is
fixed, so for random labels the cost is smallest for the prior whose values at
different inputs are uncorrelated, the squiggliest one; for labels that vary
smoothly with the input we expect the order to reverse. This is a direction
only. Beyond the inequality it is untested here: we have not measured this cost
for networks and do not know whether it predicts anything about training.

\paragraph{Conclusion.} The squiggly line picture is correct as geometry and
does not explain why overparameterized networks learn. The distance it uses is
not the one that training measures.
